\section{Introduction}
An observer does not retain a physical state merely by naming it. The observer
receives a history, chooses actions, and stores a representation from which
future decisions or predictions must be made. There are consequently three
different questions. Which future laws are determined by that history? Which
parts of those laws can a second experiment reproduce? How much persistent
state is required to reproduce them causally? The answers need not be ordered
by a single notion of information.

This paper starts with a positive causal experiment and separates these three
questions. The starting data are a preparation, transition--observation
kernels, executable future tests, admissible interventions, and a resource
interface. A posterior, a predictive quotient, and the law of attained
predictions are derived from these data. Smooth structures, information
matrices, asymptotic rates, and operator realizations are additional
structures, not axioms of the experiment.

The distinction is already visible in a binary example. If a hidden bit is
observed exactly, two persistent labels retain its posterior prediction
without loss. Adding independent Gaussian noise makes the observation less
informative in the statistical order, but the posterior then ranges
continuously between zero and one. No fixed finite number of labels preserves
all of those posterior predictions exactly. There is no violation of data
processing: the two excess losses use different full-observation oracles.
The complete calculation appears in Section~\ref{sec:boundaries}.

\subsection{The quantitative results}
Our first quantitative result concerns a prepared predictive-state model.
Let $s_t$ be its exact state and let
\[
 s_{t+1}=F_t(s_t,a_t,y_{t+1}).
\]
Suppose the same-input update is $L_t$-Lipschitz on a common admissible
domain. Cover that domain by $M_t$ representatives with radius $r_t$,
and allow an implementation error $\eta_t$ in each update. One causal
label machine then has error bounded by the deterministic recurrence
\begin{equation}\label{eq:intro-recursion}
 b_0=0,\qquad b_{t+1}=L_tb_t+r_{t+1}+\eta_{t+1}.
\end{equation}
The same machine also defines a finite-state model, provided the next-report
law is Lipschitz in the state. If its total-variation Lipschitz constant is
$\kappa_t$, the entire visible-path laws differ by at most
\begin{equation}\label{eq:intro-path-error}
 \min\left\{1,\sum_{t=0}^{N-1}\kappa_tb_t\right\},
\end{equation}
uniformly over the common feedback policy class. The policies need not be
continuous. They agree on equal observed histories, which is enough for the
kernel telescoping proof. For stopped experiments, the sum is charged only
while the reference experiment remains active. A policy that itself retains
$D_t$ states yields a controller with at most $D_tM_t$ states, not $M_t$.
This is Theorem~\ref{thm:emulation}. It distinguishes same-input prediction
accuracy from closed-loop comparison instead of inferring the latter from
the former.

The second result supplies the geometry and the actual acquisition law in a
dynamic model rather than postulating them. Fix $d\ge1$ and let the hidden
state take values in $\{0,\ldots,d\}$. Between observations it stays put with
probability $1-\gamma$ and is replaced by an independent uniform state with
probability $\gamma$, where $4/5\le\gamma\le1$. A command
$u\in[-1,1]^d$ produces a binary report $y\in\{-1,1\}$ with
\begin{equation}\label{eq:intro-sensor}
 \Pp(y\mid X=i,u)=
 \begin{cases}
  1/2,&i=0,\\
  (1+y\epsilon\tanh u_i)/2,&1\le i\le d,
 \end{cases}
 \qquad 0\le\epsilon\le1/4.
\end{equation}
At a selected checkpoint, a terminal query is drawn only after the label has
been formed. Query $i$ returns one binary outcome of mean
$1/2+\tau\1_{\{X=i\}}$, where $\tau=1/4$ and $1\le i\le d$.
All acquisition reports, both signs, are retained in the preparation law.
Only one selected terminal query is executed in a run.

Let $R_{M,t}^{\rm chk}$ be the optimal checkpoint excess squared loss under
independent uniform commands, and let $\mathfrak R_{M,N}$ be the best
maximum-checkpoint excess attainable by one causal machine with at most
$M$ labels after each acquisition. Theorem~\ref{thm:refresh} proves
\begin{equation}\label{eq:intro-refresh-law}
 c_d\epsilon^2M^{-2/d}
 \le R_{M,t}^{\rm chk}
 \le \mathfrak R_{M,N}
 \le C_d\epsilon^2M^{-2/d},\qquad 1\le t\le N.
\end{equation}
The middle inequality is understood for the optimal maximum risk, so it
holds for each fixed checkpoint. The same upper-bound machine works at every
time, and its pathwise prediction error is bounded uniformly over all
command and report sequences. Thus the maximum can be replaced by a
supremum over all positive times. The constants depend on $d$ and the fixed
terminal probe, not on $N,t,M,\epsilon$, or $\gamma$ in the displayed
ranges. At $\epsilon=0$, one label is exact.

The proof is elementary but combines obligations that should not be
suppressed. Refreshing gives an interior lower bound for the pre-observation
belief. Bayes updating is uniformly contractive after refreshing. Its entire
reachable set lies within distance $O(\epsilon)$ of the uniform belief.
The command-to-posterior Jacobian has determinant of order
$\epsilon^d$, uniformly over every previous history. It follows that the
actual posterior law has a density bounded above by $C_d\epsilon^{-d}$.
This gives a lower bound against all $M$-label encoders, including independent
public or private coding randomization. A finite cover of an invariant
belief domain gives the matching causal construction. No stationary law of
the filter and no asymptotic quantization formula is assumed.

This model has a genuinely evolving latent state. It is not the static
latent-variable, conditionally independent acquisition problem of
\cite{QiA1}. Its nonpolynomial command coordinate is convenient, not an
intrinsic novelty claim: the change $v_i=\tanh u_i$ makes the likelihood
affine in $v$. The additional conclusion relative to that manuscript's stated
finite-horizon scope is the time-uniform bound obtained here from explicit
refreshing contraction. If the terminal query is also multiplied by
$\epsilon$, its physical metric changes and the order becomes
$\epsilon^4M^{-2/d}$ (Corollary~\ref{cor:double-attenuation}).

\subsection{Objects and reductions}
A statistical comparison acts first on a family of probability laws. We use
$\Gamma:\cE\to\cF$ to mean that the source experiment $\cE$ can implement
$\cF$ by a parameter-independent causal machine. A reduction includes its
strategy lift, state, observation interface, and resource certificate.
Serial composition adds total-variation defects and composes budget maps.
For a common bounded decision loss, this gives a risk inequality. It does
not assert that two different posterior oracles have ordered coding losses.

The minimal predictive quotient is defined by equality of all designated
future tests, with equal continuation interfaces. It is an exact semantic
object. In a finite-state realization its updates are explicit; for a
general measurable quotient, standard Borel realizability is a separate
obligation. We give a compact continuous realization theorem rather than
assuming that every quotient is a smooth state space.

\subsection{Relation to established theories}
Kernel composition, disintegration, sufficiency, and their categorical
organization are established foundations; see Fritz~\cite{Fritz2020} and
Fritz--Gonda--Perrone--Rischel~\cite{FGPR2023}. Blackwell's comparison
of experiments~\cite{Blackwell1953} concerns decision performance under
garbling. Our directional simulation inequalities use that classical idea;
we retain the timing and state required to implement a garbling.
The introduction of a new project name does not make these results new.

Predictive-state and causal-state representations already define state by
future distributions, including action-conditional versions
\cite{ShaliziCrutchfield2001,LittmanSuttonSingh2002,BarnettCrutchfield2015}.
Our quotient theorem specializes to that principle. The present presentation
makes the future test class, preparation, zero-evidence convention, and
resource interface explicit, and supplies sufficient conditions for a
measurable realization.

Quantized nonlinear filtering and finite model approximation have an
extensive literature; examples include
\cite{PagesPham2005,KalogeriasPetropulu2016,SaldiYukselLinder2017}.
Approximating a signal on a finite grid while storing an arbitrary vector of
posterior weights is not the same resource as storing one of $M$ complete
prediction labels. Conversely, the cardinality bounds proved here are not
bounds on all numerical workspace or on the size of a read-only program.
The state error recurrence and the path-kernel comparison use standard
stability and telescoping arguments. The contribution of
Theorem~\ref{thm:emulation} is their explicit joint statement under feedback,
stopping, numerical defects, and simulator-state accounting. No assertion is
made that these abstract ingredients have no antecedents.

The preservation of Fisher information is local and need not imply
sufficiency, as explained by Pollard~\cite{Pollard2013}. We prove the
finite-model score formula and state the exact QMD interface. Exponential
integral--entropy duality is classical; its proof is included to fix the
risk parameter and the nonlinear pullback under a random channel.
The contact-fibre calculation of \cite{QiA2} is connected through linear
observations of positive functionals, not identified with an entire
statistical experiment.

\subsection{Reading the paper}
Sections~\ref{sec:experiments}--\ref{sec:reductions} establish the objects,
quotients, and executable reductions. Section~\ref{sec:information} fixes
the Bayesian, entropic, and score identities.
Sections~\ref{sec:checkpoint}--\ref{sec:refreshing} contain the quantitative
results and their full proofs. Sections~\ref{sec:response} and
\ref{sec:pressure} connect normalized response and exponential transforms.
Section~\ref{sec:boundaries} gives explicit limitations, and
Section~\ref{sec:interfaces} states the interfaces with subsequent geometry
and dynamical work. The appendices contain measurable and finite-experiment
details used in the proofs. There is no dependence on an unproved theorem
of a later installment.
